\documentclass[a4paper,12pt]{article}

\usepackage[utf8]{inputenc}
\usepackage[T1]{fontenc}
\usepackage[italian]{babel}
\usepackage[margin=2.5cm]{geometry}
\usepackage{graphicx}
\usepackage{grffile}
\usepackage{booktabs}
\usepackage{setspace}
\usepackage{titlesec}
\usepackage{float}
\usepackage{ifthen}
\usepackage{xcolor}
\usepackage{tcolorbox}
\usepackage{enumitem}
\usepackage[colorlinks=true,linkcolor=black,urlcolor=primaryblue,citecolor=primaryblue]{hyperref}

\definecolor{primaryblue}{RGB}{0,102,204}
\definecolor{secondaryblue}{RGB}{51,153,255}
\definecolor{lightgray}{RGB}{245,245,245}
\definecolor{darkgray}{RGB}{100,100,100}

\titleformat{\section}
  {\Large\bfseries\color{primaryblue}}
  {\thesection}{1em}{}

\setlength{\parskip}{4pt}
\setlength{\parindent}{0pt}

\setlist[itemize]{leftmargin=*,itemsep=3pt}
\setlist[enumerate]{leftmargin=*,itemsep=3pt}

\graphicspath{{./}{../assets/images/}{./images/}}

\begin{document}

\begin{center}  \IfFileExists{../../../assets/Logo.jpg}{%
    \includegraphics[width=6cm,height=3cm,keepaspectratio]{../../../assets/Logo.jpg} \\[0.8cm]
  }{%
    \fbox{\parbox[c][2.5cm][c]{6cm}{\centering Logo non trovato\\(Logo.jpg)}}\\[0.5cm]
  }
  
  {\Large\bfseries\color{primaryblue} BugBusters}\\[0.3cm]
  {\small\color{darkgray} Email: \texttt{bugbusters.unipd@gmail.com}} \\[0.1cm]
  {\small\color{darkgray} Gruppo: 4} \\[0.5cm]

  {\large\bfseries Università degli Studi di Padova}\\[0.3cm]
  {\small Laurea in Informatica}\\[0.2cm]
  {\small Corso: Ingegneria del Software}\\[0.2cm]
  {\small Anno Accademico: 2025/2026}\\[0.8cm]

  {\Huge\bfseries\color{primaryblue} Verbale Interno}\\[0.3cm]
  {\Large\color{secondaryblue} 20 ottobre 2025}\\[0.8cm]
\end{center}

\begin{center}
\begin{tcolorbox}[colback=lightgray,colframe=primaryblue,width=0.85\textwidth,arc=3mm,boxrule=0.5pt]
\begin{tabular}{@{}ll@{}}
\textbf{Redattori}    & Leonardo Salviato, Alberto Pignat \\
\textbf{Verificatore}    & Alberto Autiero \\
\textbf{Uso}          & Interno \\
\textbf{Destinatari}  & Prof. Tullio Vardanega, Prof. Riccardo Cardin \\
\end{tabular}
\end{tcolorbox}
\end{center}

\vspace{0.5cm}

\begin{center}
\begin{tcolorbox}[colback=secondaryblue!10,colframe=secondaryblue,width=0.9\textwidth,arc=3mm,boxrule=0.8pt,title={\bfseries Abstract}]
Verbale della Terza riunione avvenuta in via telematica per confronto su domande per le aziende M31 e Miriade, creazione organizzazione su github, proposta di colloquio con l'azienda Eggon, suddivisione compiti futuri, creazione prototipo sito web.
\end{tcolorbox}
\end{center}

\newpage

\tableofcontents
\newpage


\section{Informazioni generali}

\begin{itemize}
    \item \textbf{Tipo riunione:} interna
    \item \textbf{Piattaforma:} Discord
    \item \textbf{Data:} 20/10/2025
    \item \textbf{Orario di inizio:} 14:00
    \item \textbf{Orario di fine:} 16:00
    \item \textbf{Presenti:}
    \begin{itemize}[leftmargin=1.5em, itemsep=3pt, label={\rule[0.5ex]{0.4em}{0.4em}}]
        \item Alberto Autiero
        \item Alberto Pignat
        \item Marco Piro
        \item Marco Favero
        \item Linor Sadè
        \item Luca Slongo
        \item Leonardo Salviato
    \end{itemize}
    \item \textbf{Assenti:}

\end{itemize}

\section{Ordine del giorno}

\begin{enumerate}
    \item Confronto su domande per le aziende M31 e Miriade
    \item Proposta di colloquio con l'azienda Eggon
    \item Creazione organizzazione su github
    \item Suddivisione compiti futuri
    \item Creazione prototipo sito web
\end{enumerate}

\section{Svolgimento}

\begin{itemize}
    \item La riunione è iniziata alle ore 14:00.
    \item Il team ha analizzato in dettaglio la risposta pervenuta da Eggon e deciso di proporre un colloquio delucidativo con l'azienda.
    \item Poi son state discusse le domande da porre alle aziende M31 e Miriade relative ai loro capitolati
    \item L'account github creato in precedenza é stato convertito in un'organizzazione.
    \item É stato creato un prototipo di sito web per esporre ordinatamente i documenti relativi al progetto.
    \item Sono stati assegnati a ciascuno dei compiti da adempiere in vista della prossima riunione.
\end{itemize}

\section{Decisioni prese}

\begin{itemize}
    \item \textbf{Conversione in organizzazione:} É stato deciso di convertire l'account github in un'organizzazione per questioni logiche e formali.
    \item \textbf{Conferma incontro Eggon:} Son stati decisi gli orari da proporre ad Eggon per un colloquio
\end{itemize}

\section{Todo}

\begin{tcolorbox}[colback=secondaryblue!8,colframe=secondaryblue!60,arc=2mm,boxrule=0.5pt,left=10pt,right=10pt]
\begin{itemize}[topsep=5pt]
    \item Studiare un versionamento comune da usare all'interno dei documenti
    \item Redigere i paragrafi da inserire nella "Analisi dei capitolati" 
    \item Creare il documento "Analisi dei capitolati"
    \item Trovare un modo di hostare il sito web tramite github
    \item Preparare domande per un eventuale colloquio Eggon
    \item Prossima riunione: 22/10/2025 alle ore 14:00
\end{itemize}
\end{tcolorbox}

\vfill
\begin{center}
    {\small\color{darkgray} Documento redatto e approvato dal gruppo BugBusters.}
\end{center}

\end{document}